%% kapa: exam=1 round=19 year=2008 period=1 form=A subject=appraisal_law subject_ko=부동산관계법규 seq=21 number=61
%% source: Q-Net 감정평가사 기출문제 (cst00309) artlSeq=1001962 file=extracted/1008/1001962/2008년 제19회 감정평가사 1차시험 기출문제/1교시.hwp
%% conversion: kordoc 4.24.51 -> tools/gen1.py (auto); review=proofread-sonnet; flags=none
\begin{question}{61}
중앙부동산평가위원회에 관한 설명으로 옳은 것은?
\begin{choices}
  \item 위원장은 위원 중에서 호선하며, 부위원장은 위원장이 지명한다.
  \item 위원회는 위원장을 제외한 20인 이내의 위원으로 구성한다.
  \item 위원회의 회의는 국토해양부장관이 이를 소집하고, 개회 2일 전에 의안을 첨부하여 각 위원에게 통지하여야 한다.
  \item 고등교육법에 의한 대학에서 토지․주택 등에 관한 이론을 가르치는 조교수 이상의 직에 있는 자는 위원으로 위촉될 수 있다.
  \item 공무원이 아닌 위원의 임기는 2년으로 하며, 연임할 수 있다.
\end{choices}
\end{question}
\answer{4}
